\chapter{Requirements en ontwikkelproces}
\label{chap:requirements}

\section{Iteratief overleg met de klant}
De requirements lagen bij de start niet volledig vast. De opdrachtgever wist welke informatie tijdens een race nodig was, maar de precieze voorstelling en rekenregels werden duidelijker door werkende versies te bekijken. Daarom werd gedurende de zes weken veelvuldig overlegd. Een typische iteratie bestond uit een vraag of schets van de klant, een eerste implementatie, een test met live of historische data en een bespreking van onverwacht gedrag.

Deze aanpak was noodzakelijk omdat ogenschijnlijk eenvoudige velden een domeinspecifieke betekenis hebben. Zo betekent \emph{DIFF} bij GetRaceResults het tijdsverschil met de vorige wagen in de algemene rangschikking, terwijl \emph{INT} bij RIS Timing een interval kan zijn dat in seconden of in ronden wordt uitgedrukt. Een kolom met de titel \emph{IN} bleek dan weer niet voor interval te staan, maar voor de ronde waarin de beste tijd werd gereden. Zulke verschillen kunnen alleen veilig worden verwerkt wanneer zowel de website als realistische wedstrijdsituaties worden onderzocht.

\section{Evolutie van de requirements}
De eerste dashboardversie bevatte een algemene sessiekaart, een klassentabel en grafieken. Feedback leidde tot een volledig hertekend scherm volgens een schets van de opdrachtgever. De informatie werd opgesplitst in grote tijdvakken, compacte sectorvakken, vergelijkingskolommen en een pitstopbalk. Later kwamen daar een donkere modus, een apart grafiekvenster en meerdere gelijktijdige dashboards bij.

Ook de achterliggende regels evolueerden. De oorspronkelijke referentietijdlogica werd eerst verwijderd omdat de betekenis van de regel nog niet voldoende duidelijk was. Nadat de gewenste werking was verduidelijkt, werd ze opnieuw toegevoegd als configureerbare lap- en sectorreferentie met kleurstatussen. De pitstopprojectie begon als een eenvoudige aftrek van de pitstoptijd. Tests toonden vervolgens dat rekening moest worden gehouden met alle tussenliggende wagens, ook uit andere klassen, met rondeachterstanden, providerafhankelijke intervallen en FCY-omstandigheden.

De opslagstrategie veranderde eveneens. Eerst werd automatisch een nieuwe map per start gemaakt. Dat was handig, maar onveilig wanneer de applicatie per ongeluk werd gesloten: een herstart zou een nieuwe sessie beginnen. De uiteindelijke keuze is dat de gebruiker zelf één map per sessie selecteert. Bij een herstart wordt bestaande JSONL-historiek geladen en worden alleen nieuwe ronden toegevoegd.

\section{Functionele vereisten}
De gerealiseerde functionele vereisten kunnen als volgt worden samengevat:

\begin{enumerate}[label=F\arabic*.]
  \item De gebruiker configureert timing-URL, sessiemodus, opslagmap en maximaal drie autonummers.
  \item De applicatie bevraagt de timingpagina om de vijf seconden en verdeelt dezelfde data over alle dashboards.
  \item De parser ondersteunt GetRaceResults en RIS Timing en bewaart de bronprovider.
  \item Voltooide ronden worden ontdubbeld en persistent opgeslagen.
  \item Race-, practice- en qualifyingmodus tonen aangepaste analyses.
  \item Rondes onder FCY, Safety Car of andere neutralisatie blijven opgeslagen maar worden uit representatieve tempoanalyses geweerd.
  \item Inlaps en outlaps blijven opgeslagen maar tellen niet mee voor tempovergelijkingen.
  \item Groene sectoren uit een gedeeltelijk geneutraliseerde ronde kunnen afzonderlijk geldig blijven, terwijl de volledige ronde ongeldig is voor het rondetijdgemiddelde.
  \item De voorspelde rondetijd gebruikt actuele sectoren en recente sectordata van de huidige piloot.
  \item De gebruiker kan lap- en sectorreferenties aanpassen en ziet de veiligheidsmarge via kleuren en delta's.
  \item De pitstopplanner houdt rekening met openingsvensters, cooldown, verplichte stops en een instelbare pitstoptijd.
  \item Grafieken tonen piloot-, sector- en klassevergelijkingen en ondersteunen zoomen bij lange reeksen.
  \item De app kan na een crash of sluiting verdergaan vanuit dezelfde sessiemap.
\end{enumerate}

\section{Niet-functionele vereisten}
Tijdens een race wegen betrouwbaarheid en begrijpelijkheid zwaarder dan theoretische complexiteit. De berekeningen moesten daarom deterministisch, uitlegbaar en onafhankelijk van de renderer zijn. De renderer krijgt kant-en-klare viewmodellen en formatteert of berekent geen racebeslissingen zelf. Elke belangrijke domeinfunctie moest met synthetische en historische data testbaar zijn.

Verder moest de applicatie langdurig kunnen draaien zonder onbegrensde geheugengroei of overmatige opslag. De pollingfrequentie van vijf seconden levert tijdens 24 uur 17.280 meetmomenten op, maar alleen gewijzigde snapshots en voltooide rondes worden als relevante historiek bewaard. Het gebruik van één collector voor maximaal drie dashboards voorkomt onnodige netwerkbelasting.

\section{Planning over zes weken}
De uitvoering verliep niet als zes strikt gescheiden fasen. In de eerste periode lag de nadruk op parser, live collector en opslag. Daarna volgden dashboardanalyse, tests en de nieuwe UI. Pitstoplogica, voorspelling, sessiemodi en grafieken werden in opeenvolgende iteraties toegevoegd. In de laatste fase kwamen provideredgecases, FCY-projecties, distributie, automatische updates en herstel na herstart aan bod.

Deze iteratieve volgorde week af van een klassiek vooraf vastgelegd stageplan, maar paste beter bij de beschikbaarheid van live demo's en klantfeedback. Bugs die alleen in een bijna volledige race zichtbaar werden, hadden een grotere impact op de prioriteiten dan oorspronkelijk geplande uitbreidingen.

